\scanpage{101}
qui pour $S$ variable est compatible \unsure{aux|avec} chgt de base, et donne donc un morphisme de multiplicités modulaires \inserted{(si $G,G'$ sont stables)}
\[
(132)\qquad M_G\longrightarrow M_{G'},
\]
avec fonctorialité évidente pour les composés. Ainsi $M_G$ est un \uline{foncteur contravariant en} $G$ pour les morphismes de normalisation \inserted{de m.graphes stables}. Cette fois, on a

\uline{Proposition} a) Si $G\longrightarrow G'$ est un morphisme de normalisation \uline{strict}, alors (132) est un \uline{iso.} de multiplicités

b) Soient $G_\alpha$ des \deleted{\unsure{composantes connexes|graphes connexes}}\inserted{m.graphes stables} \deleted{\unsure{de}} $G$ le m.graphe somme des $G_\alpha$. Alors\unsure{:|,} le morphisme
\[
(133)\qquad M_G\unsure{\rightharpoondown|{-}}\prod_\alpha, M_{G_\alpha}
\]
\unsure{est} un iso. de multiplicités.
\margin{\uline{NB} On remplace la notation $\Pi M_G$ par $\Pi_1M_G$, pour ne pas \unsure{confondre} $\Pi_1\deleted{\unsure{\Pi|M_G}}$ le signe du produit}

\uline{Cor} $\Pi_1M_G$ est un foncteur contravariant en \deleted{\unsure{le}} m.graphe stable $G_{\deleted{\unsure{\alpha|f}}}$ pour les
\margin{100}
